
Las instalaciones de SKYFILMS destinadas a las actividades de mantenimiento, inspección, almacenamiento y control técnico de los Sistemas de Aeronaves No Tripuladas (UAS) están adecuadamente equipadas y organizadas para garantizar el cumplimiento de los procedimientos establecidos en este Manual de Control de Mantenimiento (MCM) y en los Reglamentos Aeronáuticos de Colombia (RAC 100).

Estas instalaciones permiten ejecutar de manera segura y eficiente las tareas de mantenimiento preventivo, correctivo, limpieza, calibración de equipos, carga de baterías y almacenamiento de componentes, asegurando condiciones que preservan la integridad técnica de las aeronaves y la seguridad del personal.

\section{Características}
Las labores de mantenimiento se realizan en las instalaciones principales de SKYFILMS, ubicadas en Ibagué, Tolima, las cuales cuentan con un área delimitada y señalizada para la manipulación y mantenimiento de aeronaves no tripuladas.

Las áreas están organizadas de acuerdo con su función y cuentan con los elementos mínimos requeridos por la Aerocivil, incluyendo:

\begin{itemize}
    \item \textbf{Área de mantenimiento técnico:} espacio destinado al mantenimiento menor, mayor y correctivo, inspección general y verificación de sistemas de los UAS.
    
    \item \textbf{Área de limpieza y preparación:} destinada a la limpieza superficial y acondicionamiento previo y posterior al vuelo.
    
    \item \textbf{Área de carga y gestión de baterías:} equipada con cargadores, sistemas de ventilación y elementos de seguridad contra incendios, conforme a los lineamientos de seguridad eléctrica.

    \item \textbf{Zona de almacenamiento técnico:} destinada al almacenamiento de aeronaves, herramientas, repuestos, equipos de control remoto y sensores.


\end{itemize}

En caso de requerirse mantenimiento especializado o intervenciones mayores, SKYFILMS podrá utilizar instalaciones de terceros aceptados, previa verificación de su idoneidad técnica y cumplimiento con la normativa aplicable.\\


\section{Estándares de Limpieza y Conservación}

\subsection{Alcance}
El presente apartado establece los lineamientos obligatorios para garantizar que todas las instalaciones utilizadas por SKYFILMS cumplan con los más altos estándares de limpieza, orden y conservación. Estas directrices aplican a las áreas destinadas al mantenimiento (menor, mayor y correctivo), almacenamiento, ensamblaje, pruebas y demás instalaciones aplicables a la organización. El objetivo principal es prevenir daños por objetos extraños, proteger la integridad de los componentes electrónicos y estructurales de los Sistemas de Aeronaves No Tripuladas (UAS).

\subsection{Estándares Específicos por Área de Trabajo}
Para asegurar la aeronavegabilidad continua de la flota, las instalaciones de SKYFILMS (propias o tercerizadas aceptadas) deben cumplir con las condiciones detalladas en la siguiente tabla de acuerdo con la naturaleza de la actividad:

\begin{longtable}{|p{0.22\textwidth}|p{0.34\textwidth}|p{0.34\textwidth}|}
% --- ENCABEZADO DE LA PRIMERA PÁGINA ---
\hline
\rowcolor[HTML]{EFEFEF} 
\textbf{Área de Instalación} & \textbf{Estándares de Limpieza} & \textbf{Estándares de Conservación e Infraestructura} \\ \hline
\endfirsthead

% --- ENCABEZADO PARA LAS PÁGINAS SIGUIENTES ---
\hline
\rowcolor[HTML]{EFEFEF} 
\textbf{Área de Instalación} & \textbf{Estándares de Limpieza} & \textbf{Estándares de Conservación e Infraestructura} \\ \hline
\endhead

% --- PIE DE PÁGINA CUANDO LA TABLA CONTINÚA ---
\hline
\multicolumn{3}{r}{\textit{Continúa en la siguiente página...}} \\ 
\endfoot

% --- PIE DE PÁGINA FINAL CON EL RÓTULO DEBAJO ---
\hline
\caption{Estándares Específicos por Área de Trabajo} \label{tab:estandares_limpieza} \\
\endlastfoot

% --- CONTENIDO DE LA TABLA ---

\vspace{2pt} \raggedright \textbf{Área de Mantenimiento (Menor, Mayor y Correctivo)} & 
\begin{minipage}[t]{\linewidth}
\vspace{2pt}
\raggedright
\begin{itemize}[leftmargin=12pt, topsep=0pt, partopsep=0pt, parsep=0pt, itemsep=4pt]
    \item Superficies de trabajo (mesas) libres de polvo, grasas, líquidos y residuos de materiales antes de iniciar cualquier tarea.
    \item Prohibición estricta de consumo de alimentos o bebidas en las estaciones de trabajo.
    \item Limpieza y desinfección de herramientas manuales al finalizar cada jornada.
\end{itemize}
\vspace{2pt}
\end{minipage} & 
\begin{minipage}[t]{\linewidth}
\vspace{2pt}
\raggedright
\begin{itemize}[leftmargin=12pt, topsep=0pt, partopsep=0pt, parsep=0pt, itemsep=4pt]
    \item Mesas equipadas con tapetes antiestáticos (ESD) conectados a tierra.
    \item Iluminación blanca, uniforme y sin sombras que permita la inspección visual detallada de componentes.
    \item Herramientas organizadas mediante tableros de sombras (Shadow Boards) para control de inventario y prevención de FOD.
\end{itemize}
\vspace{2pt}
\end{minipage} \\ \hline

\vspace{2pt} \raggedright \textbf{Área de Almacenamiento (Equipos, Repuestos y Consumibles)} & 
\begin{minipage}[t]{\linewidth}
\vspace{2pt}
\raggedright
\begin{itemize}[leftmargin=12pt, topsep=0pt, partopsep=0pt, parsep=0pt, itemsep=4pt]
    \item Pisos y estanterías aspirados o barridos con métodos que no levanten polvo en suspensión.
    \item Eliminación inmediata de empaques vacíos y material de desecho.
\end{itemize}
\vspace{2pt}
\end{minipage} & 
\begin{minipage}[t]{\linewidth}
\vspace{2pt}
\raggedright
\begin{itemize}[leftmargin=12pt, topsep=0pt, partopsep=0pt, parsep=0pt, itemsep=4pt]

    \item Repuestos almacenados en sus empaques originales, clasificados y etiquetados.
    \item \textbf{Baterías (Li-Po/Li-Ion):} Almacenadas en gabinetes ignífugos, a temperatura controlada (15°C - 25°C) y con niveles de carga de almacenamiento (storage voltage) según manual del fabricante.
\end{itemize}
\vspace{2pt}
\end{minipage} \\ \hline

\vspace{2pt} \raggedright \textbf{Área de Ensamblaje} & 
\begin{minipage}[t]{\linewidth}
\vspace{2pt}
\raggedright
\begin{itemize}[leftmargin=12pt, topsep=0pt, partopsep=0pt, parsep=0pt, itemsep=4pt]
    \item Entorno estéril libre de partículas en el aire que puedan ingresar a motores o placas de circuitos.
    \item Segregación de desechos generados durante el ensamblaje (cintas, cables cortados, precintos).
\end{itemize}
\vspace{2pt}
\end{minipage} & 
\begin{minipage}[t]{\linewidth}
\vspace{2pt}
\raggedright
\begin{itemize}[leftmargin=12pt, topsep=0pt, partopsep=0pt, parsep=0pt, itemsep=4pt]
    \item Superficies niveladas y acolchadas para evitar rayones o daños estructurales en el fuselaje o brazos de los UAS durante el armado.
    \item Acceso restringido a personal no autorizado para evitar manipulación indebida o accidentes.
\end{itemize}
\vspace{2pt}
\end{minipage} \\ \hline

\vspace{2pt} \raggedright \textbf{Área de Pruebas (Indoor/Outdoor)} & 
\begin{minipage}[t]{\linewidth}
\vspace{2pt}
\raggedright
\begin{itemize}[leftmargin=12pt, topsep=0pt, partopsep=0pt, parsep=0pt, itemsep=4pt]
    \item Implementación del programa de control FOD: inspección visual previa para retirar piedras, tornillos sueltos, plásticos o cualquier objeto que pueda ser succionado por los rotores.
\end{itemize}
\vspace{2pt}
\end{minipage} & 
\begin{minipage}[t]{\linewidth}
\vspace{2pt}
\raggedright
\begin{itemize}[leftmargin=12pt, topsep=0pt, partopsep=0pt, parsep=0pt, itemsep=4pt]
    \item Demarcación clara de las zonas de encendido y pruebas de motores.
    \item Disponibilidad de extintores (Clase ABC y D para baterías) y botiquín de primeros auxilios en puntos visibles y de fácil acceso.
    \item Protección contra factores climáticos directos si las pruebas se realizan en recintos cerrados o semi-abiertos.
\end{itemize}
\vspace{2pt}
\end{minipage} \\ \hline

\end{longtable}

\subsection{Frecuencia de Aseo y Mantenimiento de Instalaciones}

\begin{itemize}
    \item \textbf{Rutina Diaria (Al cierre de jornada):} Organización de herramientas, recolección de residuos ordinarios y limpieza de superficies de trabajo.
    
    \item \textbf{Rutina Semanal:} Limpieza profunda de pisos, estanterías de repuestos y revisión de los niveles de almacenamiento de las baterías.

    \item \textbf{Rutina Mensual:} Revision de instalaciones por parte del Gerente SMS para verificar iluminación, integridad de los tapetes antiestáticos, estado de extintores y cumplimiento general de estos estándares.
\end{itemize}


